\documentclass[11pt]{article}

\usepackage[a4paper,margin=1in]{geometry}
\usepackage{amsmath}
\usepackage{amssymb}
\usepackage{booktabs}
\usepackage{hyperref}
\usepackage{xcolor}
\usepackage{listings}

\lstset{
  basicstyle=\ttfamily\small,
  breaklines=true,
  frame=single,
  columns=fullflexible
}

\title{Robot + Object No-Object-Euler World Model:\\Coupling-driven Object Dynamics with an Online Latent Context Filter}
\author{}
\date{}

\begin{document}
\maketitle

\section{Overview}

This document describes the architecture implemented under
\begin{center}
\texttt{scripts/world\_model/Physics/Robot\_and\_Object/No\_object\_Euler}
\end{center}
which is a fork of the robot+object MCGDF model with two deliberate
structural changes:
\begin{enumerate}
  \item \textbf{No Newton--Euler integrator on the object side.}  The
    object's twist is produced by blending the analytic gripper twist
    $J(q)\dot q$ into the previous object twist via a contact-gated
    coupling factor $\alpha_t \in [0,1]$.  The dataset's ground-truth
    physical context $\xi = [m_o,\ I_b,\ \mu_o]$ no longer appears
    anywhere in the model's forward pass; it is used only by
    training-time supervised losses and InfoNCE.
  \item \textbf{Online latent context update.}  The variational latent
    $z$ is no longer a single one-shot encoding of a fixed history
    window: it is refined step by step from the per-step object-state
    prediction residual using a sigmoid-gated innovation cell.
\end{enumerate}
The robot side (DeLaN $H/c/g$, the residual head $r_\theta$, joint
damping/friction, and the Newton-third-law coupling
$\tau_{\mathrm{contact}} = -J_c^\top F_c$) is inherited from the MCGDF
model unchanged.

This design follows the project proposal's Pi-CDD direction:
``Action--Reaction'' inference of an evolving per-step context from
observable contact data, without consuming GT physical labels at
deployment.  See
\texttt{../MCGDF/robot\_object\_mcgdf\_description.tex},
Section ``Coupling-driven object dynamics with an online latent
context update'' for the planning discussion that motivates the
present implementation.

\section{State, Input, and Context}

The state, the per-step input, and the per-episode context are
structurally the same as in MCGDF:
\begin{equation}
  s_t =
  \begin{bmatrix}
    q_t \\ \dot q_t \\ p^o_t \\ r^o_t \\ v^o_t \\ \omega^o_t
  \end{bmatrix} \in \mathbb{R}^{18 + 13} = \mathbb{R}^{31},
  \quad
  \tau_t \in \mathbb{R}^9,
  \quad
  \xi = \bigl[\,m_o,\ I_b\,\mathrm{(9)},\ \mu_o\,(3)\,\bigr] \in \mathbb{R}^{13}.
\end{equation}
\paragraph{Asymmetry with MCGDF.}  In MCGDF, $\xi$ enters the
\texttt{ContactWrenchHead} input concatenation and the Newton--Euler
step itself.  In the No-Object-Euler architecture both of those
consumers are removed:
\begin{itemize}
  \item The Newton--Euler integrator no longer exists
    (Section~\ref{sec:object-side}).
  \item The contact head is built with \texttt{omit\_object\_context=True}
    (Section~\ref{sec:contact-head}), so $\xi$ is never concatenated.
\end{itemize}
$\xi$ is still loaded by the dataset and supplied to the model
constructor, but it is consumed only by training-time auxiliary terms
(supervised object-dynamics targets, the InfoNCE soft-similarity
kernel, and the residual baseline computation).

\section{Forward Dynamics}

\subsection{Robot side (unchanged)}

The robot inverse-dynamics equation is identical to MCGDF:
\begin{equation}
  H(q)\,\ddot q
  = \tau + r_\theta(q,\dot q,\tau)
  - c(q,\dot q) - g(q)
  - d\dot q - f\tanh(\dot q/\varepsilon_f)
  - J_c(q)^\top F_c(s,\tau,z),
  \label{eq:noe-robot}
\end{equation}
where $H$, $c$, $g$ come from the DeLaN block and $r_\theta$, $F_c$
come from learned heads.  The crucial difference from MCGDF is the
argument of the contact head: $F_c$ now depends on $(s,\tau,z)$ only,
not on $\xi$.

The damping double-counting caveat carries over from MCGDF; the
default is \texttt{--omit\_damping} so the recorded
\texttt{applied\_torque} (which already contains the implicit-PD
damping torque) is used as-is.

\subsection{Object side: coupling-based dynamics}
\label{sec:object-side}

In place of
$m_o\dot v^o = m_o g + F_c^{\mathrm{lin}}$ and the gyroscopic Euler
equation, the object substate is integrated by the following block.
Let $p_g(q)$ be the analytic gripper-tip position
(\texttt{FrankaForwardKinematics}), and let
$(v_g, \omega_g) = J(q)\,\dot q$ be the gripper twist in world frame
(the geometric Jacobian product, also analytic).  Define
\begin{align}
  \mathbf{d}_t &= p_g(q_t) - p^o_t, \label{eq:noe-displacement}\\
  \alpha_t &= \mathrm{sigmoid}\!\bigl(\mathrm{MLP}_\alpha\bigl(\|F_c^{\mathrm{lin}}\|,\,\mathbf{d}_t,\,z_t\bigr)\bigr) \;\in\;[0,1]. \label{eq:noe-alpha}
\end{align}
$\alpha_t$ has \texttt{coupling\_components} channels (default 1; set
to 2 to enable separate $(\alpha^{\mathrm{rigid}},
\alpha^{\mathrm{slide}})$).  The object twist is then advanced by
blending the gripper twist into the previous object twist, with small
learned passive decays $\rho_v, \rho_\omega \in (0,1)$
parameterised via a logit so they stay bounded:
\begin{align}
  v^o_{t+1}      &= \alpha_t\,v_g + (1-\alpha_t)\,\rho_v\,v^o_t,      \label{eq:noe-vobj}\\
  \omega^o_{t+1} &= \alpha_t\,\omega_g + (1-\alpha_t)\,\rho_\omega\,\omega^o_t. \label{eq:noe-womega}
\end{align}
Position is advanced by semi-implicit Euler, and orientation by the
closed-form exponential map (both unchanged from MCGDF):
\begin{align}
  p^o_{t+1} &= p^o_t + v^o_{t+1}\,\Delta t, \label{eq:noe-pobj}\\
  r^o_{t+1} &= r^o_t \otimes \exp\!\bigl(\tfrac{1}{2}\omega^o_{t+1}\,\Delta t\bigr). \label{eq:noe-robj}
\end{align}
None of $m_o,I_b,\mu_o$ appears anywhere in
Eqs.~\eqref{eq:noe-displacement}--\eqref{eq:noe-robj}.

\paragraph{What is gained and lost.}
\emph{Gained.} The object-side parameter budget collapses to one small
\texttt{CouplingHead} MLP plus two logit scalars
$(\rho_v, \rho_\omega)$.  There is no mass-matrix solve, no
gyroscopic cross product, no rotation-into-world-frame inertia, and
no positive-definite-inertia parameterisation.  The integrator
forward pass is finite by construction (everything is a scalar blend
of finite inputs); the original MCGDF blowup mode through
$\omega^o \to \infty$ in the quaternion exp map cannot occur here.\\
\emph{Lost.} The object side no longer satisfies $F = m a$ or the
gyroscopic Euler equation.  Motions the new integrator cannot
represent include ballistic flight (a thrown cube), free spinning
after release, and any cube motion driven by an impulse other than
the gripper.  The 6-D wrench $F_c$ now enters the object integrator
only through its magnitude $\|F_c^{\mathrm{lin}}\|$ as a gating
feature, not as a force.  The robot-side action--reaction coupling
$-J_c^\top F_c$ is preserved.

\section{Online Latent Context Update}
\label{sec:latent-update}

The latent $z$ is no longer the encoder mean held fixed across the
whole rollout.  It is updated step by step from the per-step
prediction residual using a small sigmoid-gated cell
(\texttt{LatentInnovationUpdate}).  Let
$\mathrm{Proj}_{\mathrm{obj}}: \mathbb{R}^{31} \to \mathbb{R}^{13}$
select the object substate $(p^o, r^o, v^o, \omega^o)$ from a full
state vector.  Given the prediction $\hat s_t$ produced under
$z_{t-1}$ from a one-step call to
\texttt{RobotObjectMCGDFStep.forward}, define
\begin{equation}
  e_t \;=\; \mathrm{Proj}_{\mathrm{obj}}\!\bigl(s_t^{\mathrm{obs}} - \hat s_t\bigr) \in \mathbb{R}^{13}.
  \label{eq:noe-residual}
\end{equation}
The latent is then updated by
\begin{equation}
  \begin{aligned}
  [\boldsymbol{g}_t,\ \boldsymbol{\delta}_t] &= \mathrm{MLP}_{\mathrm{LU}}\!\bigl(z_{t-1},\,e_t,\,F_c^{t-1}\bigr) \in \mathbb{R}^{2\,\mathrm{latent\_dim}}, \\
  z_t &= z_{t-1} \;+\; \mathrm{sigmoid}(\boldsymbol{g}_t)\;\odot\;\tanh(\boldsymbol{\delta}_t).
  \end{aligned}
  \label{eq:noe-latent-update}
\end{equation}
$\mathrm{sigmoid}(\boldsymbol{g}_t)$ is the gate that controls update
strength; $\tanh(\boldsymbol{\delta}_t)$ is the bounded delta
direction.  The update is a Kalman-innovation analog: the only way
$z$ can change is to follow the surprise $e_t$, weighted by the
contact wrench evidence.

\paragraph{Where the update runs (filter window).}
\texttt{filter\_steps} controls how many closed-loop steps run before
the open-loop rollout.  At training time, given a history of
$K{+}1$ observed states (and the corresponding torque sequence), the
filter consumes the last $L_{\mathrm{filter}}$ observed transitions:
\begin{lstlisting}
z = ContextEncoder(history)                  # prior from warm-up window
for k in range(K - L_filter, K):             # closed-loop filter window
    s_hat = OneStep(history[k], tau[k], z)
    e     = Proj_obj(history[k+1] - s_hat)
    z     = z + sigmoid(gate) * tanh(delta)  # Eq. (noe-latent-update)
for h in range(rollout_horizon):             # open-loop rollout window
    s = OneStep(s, future_torques[h], z)     # z is frozen here
\end{lstlisting}
At evaluation time the same code runs.  Crucially the filter window
reads only $s_t^{\mathrm{obs}}$ (joint encoders and the recorded
cube pose), never the GT $\xi$, so it is deployment-fair.

\paragraph{What gets backpropagated.}
The supervised losses act on the open-loop rollout predictions.
Their gradients flow backward through the rollout, through
$z_{K{+}L_{\mathrm{filter}}}$, and through the
$L_{\mathrm{filter}}$ steps of
Eq.~\eqref{eq:noe-latent-update} into the LatentInnovationUpdate
cell, the encoder, and (via the per-step closed-loop predictions)
into the dynamics heads.  Truncated BPTT over the filter window
keeps the cost affordable when $L_{\mathrm{filter}}$ is set large.

\section{Neural Network Architecture}
\label{sec:architecture}

\subsection{Sub-network inventory}

\begin{table}[h]
\centering
\small
\begin{tabular}{p{0.30\linewidth} p{0.30\linewidth} p{0.18\linewidth} p{0.10\linewidth}}
\toprule
Sub-network (class) & Input & Output & Hidden \\
\midrule
\texttt{g\_net} (MLP)                & $q\in\mathbb{R}^9$        & $g\in\mathbb{R}^9$         & 256 \\
\texttt{l\_diag\_net} (Deriv.\,MLP)  & $q\in\mathbb{R}^9$        & $\ell_d\in\mathbb{R}^9$    & 256 \\
\texttt{l\_offdiag\_net} (Deriv.\,MLP) & $q\in\mathbb{R}^9$      & $\ell_o\in\mathbb{R}^{36}$ & 256 \\
\texttt{ResidualTorqueHead}          & $(q,\dot q,\tau)\in\mathbb{R}^{27}$ & $r_\theta\in\mathbb{R}^9$ & 128 \\
\texttt{ContactWrenchHead}           & $(s,\tau)\in\mathbb{R}^{40}$ ($\xi$ omitted) & $F_c\in\mathbb{R}^6$ & 128 \\
\texttt{CouplingHead}                & $(\|F_c^{\mathrm{lin}}\|,\mathbf{d}_t,z_t)\in\mathbb{R}^{4+\mathrm{ld}}$ & $\alpha_t\in[0,1]^{c}$ & 64 \\
\texttt{LatentInnovationUpdate}      & $(z,e_t,F_c)\in\mathbb{R}^{\mathrm{ld}+19}$ & $z_{t}\in\mathbb{R}^{\mathrm{ld}}$ & 64 \\
\texttt{ContextEncoder}              & flat history seq.\ ($K{+}1)\times(\dim s + \dim\tau)$ & $(\mu,\log\sigma^2)\in\mathbb{R}^{2\cdot\mathrm{ld}}$ & 256 \\
\texttt{JointDampingFrictionParams} (opt.) & --- & $9{+}9$ scalars & --- \\
$(\rho_v,\rho_\omega)$ logit scalars & --- & $2$ scalars & --- \\
\bottomrule
\end{tabular}
\caption{Trainable sub-networks of the No-Object-Euler model.
$\mathrm{ld}=\texttt{latent\_dim}$; $c=\texttt{coupling\_components}\in\{1,2\}$.}
\label{tab:noe-subnets}
\end{table}

\paragraph{What is \emph{not} learned.}
Identical to MCGDF on the robot side: Franka FK and Jacobian,
quaternion exponential map, and semi-implicit Euler integration.  On
the object side: the kinematic blend
Eqs.~\eqref{eq:noe-vobj}--\eqref{eq:noe-robj} and the
semi-implicit Euler position update.

\subsection{End-to-end data flow per step}

\begin{lstlisting}
state s = [q (9), qdot (9), p_o (3), r_o (4), v_o (3), omega_o (3)]  = 31
input  tau (9)
ctx    xi  = [m_o (1), I_b (9), mu_o (3)]      # supervised-only, not consumed

Structural rigid-body (DeLaN):
  g_net(q)         -> g (9)
  l_diag_net(q)    -> ell_d (9)
  l_offdiag_net(q) -> ell_o (36)  -> H = LL^T + eps I, c via DeLaN identity

Implicit-PD bias:
  r_theta(q, qdot, tau) -> 9

Forward kinematics (analytic, fixed):
  FK(q) -> p_ee (3), R_ee (3x3)
  J(q)  -> 6 x 9, with last 2 finger columns zero

Contact wrench (NO xi):
  F_c(s, tau, z)   -> 6
  tau_contact      = -J(q)^T F_c       -> 9

Robot side (Newton):
  H qdd = tau + r_theta - c - g - d qdot - f tanh(qdot/eps) + tau_contact
  qdd   = solve(H, RHS)
  q_next, qdot_next via semi-implicit Euler

Coupling-based object side:
  gripper twist  (v_g, omega_g) = J(q) qdot
  displacement   d_t = p_ee - p_o
  alpha_t        = sigmoid(MLP(||F_c^lin||, d_t, z))
  v_o_next       = alpha_t * v_g + (1 - alpha_t) * rho_v   * v_o
  omega_o_next   = alpha_t * omega_g + (1 - alpha_t) * rho_w * omega_o
  p_o_next       = p_o + v_o_next * dt          (semi-implicit Euler)
  r_o_next       = r_o (x) exp(0.5 omega_o_next dt)
\end{lstlisting}

\section{Trajectory Prediction: filter window + open-loop rollout}
\label{sec:trajectory-prediction}

A predicted trajectory over a horizon $H$ is now produced in two
phases (see also Section~\ref{sec:latent-update}):

\begin{enumerate}
  \item \textbf{Closed-loop filter window} of length
    $L_{\mathrm{filter}}$.  For each transition
    $(s_{k-1}^{\mathrm{obs}}, \tau_{k-1}, s_k^{\mathrm{obs}})$
    inside the filter window, the model runs a one-step prediction
    with the current $z$ and updates $z$ by
    Eq.~\eqref{eq:noe-latent-update}.  The filter is the only place
    a model-prediction error is used as a signal for $z$.
  \item \textbf{Open-loop rollout} of length $H$.  Starting from
    $s_t = s_{t}^{\mathrm{obs}}$ at the end of the filter window, the
    model iterates
    \begin{equation}
      s_{t+h+1} = \mathrm{RobotObjectMCGDFStep}\bigl(s_{t+h},\,\tau_{t+h+1},\,z_K\bigr),
    \end{equation}
    where $z_K$ is the filter's final value held fixed.
\end{enumerate}

The reported open-loop horizon-$H$ error is therefore a function of
$(K, L_{\mathrm{filter}}, H)$.  A natural diagnostic plot is a sweep
over $L_{\mathrm{filter}}$ at fixed $(K, H)$: open-loop error should
drop monotonically with $L_{\mathrm{filter}}$ until the latent
saturates -- this is what demonstrates the value of the online
filter.

\section{Forward Kinematics and Jacobian}

Identical to MCGDF.  \texttt{FrankaForwardKinematics} encodes the
seven-link Franka chain as fixed buffers; \texttt{fk(q)} returns
$(p_{\mathrm{ee}}, R_{\mathrm{ee}})$ and \texttt{fk.jacobian(q)}
returns the geometric Jacobian $J(q) \in \mathbb{R}^{6\times 9}$ with
zero last two columns (finger DoFs).  In the No-Object-Euler model
the Jacobian is now consumed in \emph{two} places:
\begin{itemize}
  \item Robot side: $\tau_{\mathrm{contact}} = -J(q)^\top F_c$
    (Newton's third law, unchanged from MCGDF).
  \item Object side: $(v_g, \omega_g) = J(q)\dot q$, the gripper
    twist that feeds the coupling blend
    Eqs.~\eqref{eq:noe-vobj}--\eqref{eq:noe-womega}.
\end{itemize}

\section{Contact Wrench Head}
\label{sec:contact-head}

\texttt{ContactWrenchHead}'s mathematical role is the same as in
MCGDF: predict a 6-D wrench $F_c = (F_c^{\mathrm{lin}},
F_c^{\mathrm{ang}})$ in world frame.  Two API changes:
\begin{itemize}
  \item \textbf{No $\xi$ input.}  Built with
    \texttt{context\_dim=0} via the new constructor flag
    \texttt{omit\_object\_context=True} (default in this folder),
    the head consumes only $(s, \tau, z)$.  At call time, passing
    \texttt{context=None} (or any tensor with last dim 0) is
    accepted; the head silently drops $\xi$ from the concatenation.
  \item \textbf{Two consumers downstream, not one.}  The output
    drives both $\tau_{\mathrm{contact}} = -J^\top F_c$ on the robot
    side and (through its linear-component magnitude) the coupling
    gate $\alpha_t$ on the object side.  Keeping $F_c$ at 6-D
    preserves the robot-side action--reaction even though only
    $\|F_c^{\mathrm{lin}}\|$ enters $\alpha_t$.
\end{itemize}

\section{Residual Head and Pre-Contact Baseline}

Identical to MCGDF.  \texttt{ResidualTorqueHead} takes $(q,\dot q,
\tau)$ and outputs an implicit-PD bias correction $r_\theta$.  The
per-episode pre-contact baseline subtraction (Approach A of
mismatch~\#2 in the verification report) is unchanged; the trainer
shifts the per-step implicit residual target by the mean over the
pre-contact window so that contact-period excess is left for
$-J^\top F_c$ to absorb.

\section{Coupling Head and Passive Decays}
\label{sec:coupling-head}

\texttt{CouplingHead} is a small MLP with input
$(\|F_c^{\mathrm{lin}}\|, \mathbf{d}_t, z_t) \in \mathbb{R}^{4+\mathrm{ld}}$
and output passed through a sigmoid:
\begin{equation}
  \alpha_t = \mathrm{sigmoid}\bigl(\mathrm{MLP}_\alpha(\cdots)\bigr) \in [0,1]^{c},
\end{equation}
where $c = \texttt{coupling\_components}$.  When $c=1$ (default), a
single scalar gates both linear and angular blending.  When $c=2$,
$\alpha[0]$ gates the linear coupling and $\alpha[1]$ gates the
angular, allowing the model to express ``rigidly held'' vs.\
``sliding under push'' regimes asymmetrically.

The passive decays $\rho_v, \rho_\omega \in (0,1)$ are stored as raw
logits and read out as $\mathrm{sigmoid}(\cdot)$ at call time, so
they stay bounded:
\begin{equation}
  \rho_v = \mathrm{sigmoid}(\rho_v^{\mathrm{raw}}),
  \quad
  \rho_\omega = \mathrm{sigmoid}(\rho_\omega^{\mathrm{raw}}),
  \quad
  \rho_v^{\mathrm{raw}}, \rho_\omega^{\mathrm{raw}} \in \mathbb{R}.
\end{equation}
Their initial values default to $0.95$ (slight dissipation) via the
flags \texttt{--init\_rho\_v} and \texttt{--init\_rho\_omega}.

\section{Latent Innovation Update (cell details)}

\texttt{LatentInnovationUpdate} is a single hidden-layer MLP (depth
configurable) whose output is split into a sigmoid gate and a
$\tanh$ delta, both of width \texttt{latent\_dim}.  Implementation:
\begin{lstlisting}
class LatentInnovationUpdate(nn.Module):
    def __init__(self, latent_dim, residual_dim=13, contact_dim=6,
                 hidden_dim=64, depth=1):
        ...
        # in: cat([z_prev (ld), e_obj (13), F_c (6)])  =>  ld + 19
        self.backbone = MLP(ld+19, hidden_dim, depth) with SiLU
        self.head     = Linear(hidden_dim, 2*ld)

    def forward(self, z_prev, residual_obj, F_c):
        h    = self.backbone(cat([z_prev, residual_obj, F_c]))
        gate = sigmoid(self.head(h)[:, :ld])
        delta= tanh   (self.head(h)[:, ld:])
        return z_prev + gate * delta
\end{lstlisting}
\paragraph{Why object-only residual.}  We supply only
$\mathrm{Proj}_{\mathrm{obj}}(s_t^{\mathrm{obs}} - \hat s_t) \in
\mathbb{R}^{13}$ to the cell.  The robot-side residual is dominated
by joint-tracking noise once $r_\theta$ has absorbed the implicit-PD
bias, and would mostly serve as a distractor for $z$.  This is the
``object-only $e_t$'' choice from the design notes; it is enforced
inside \texttt{MultiStepRobotObjectMCGDFWorldModel.step\_filter}.

\section{Loss}

The loss menu is unchanged in structure from MCGDF -- same names,
same weights, same scaffolding -- but several terms now have
slightly different semantics under the new dynamics:
\begin{itemize}
  \item $\mathcal{L}_{\mathrm{rollout}}$: weighted MSE on robot and
    object substates.  The object-pos and quaternion components now
    measure how well the coupling-driven integrator tracks the
    recorded cube trajectory; the velocity components measure the
    blended velocity directly.
  \item $\mathcal{L}_{\mathrm{robotGT}}$: identical to MCGDF.
  \item $\mathcal{L}_{\mathrm{objectGT}}$: weights $\dot v^o, \dot\omega^o$
    against the recorded accelerations.  In the new model
    \texttt{aux["object\_lin\_acc"]} and
    \texttt{aux["object\_ang\_acc"]} are finite differences of the
    \emph{blended} velocities, so this loss now constrains the
    coupling and decay scalars to produce kinematically consistent
    velocity changes.  Practitioners may want to reduce
    \texttt{--lambda\_object\_lin\_acc} / \texttt{--lambda\_object\_ang\_acc}
    relative to MCGDF since they no longer have a Newton--Euler
    direct-acceleration interpretation; the rollout MSE on velocities
    alone is often sufficient.
  \item $\mathcal{L}_{\mathrm{DF}}, \mathcal{L}_{\mathrm{res,sup}}$:
    unchanged from MCGDF.
  \item $\mathcal{L}_{\mathrm{reg}}$: L2 on coriolis, $r_\theta$,
    and $F_c$ magnitudes -- unchanged.
  \item $\mathcal{L}_{\mathrm{KL}}$ on $z_0$: unchanged; the prior
    is N(0, I) and only the encoder-mean of the warm-up window is
    regularized.  The filter-time latents $z_K, \dots$ are not
    KL-penalized.
  \item $\mathcal{L}_{\mathrm{NCE}}$: unchanged.  The InfoNCE kernel
    still uses GT $\xi$ as the soft anchor for similarity; this is a
    training-time signal and does not enter the model's forward
    pass.
  \item $\mathcal{L}_{\mathrm{ctx,reg}}$: still available but
    deprecated for this architecture.  Setting
    \texttt{--lambda\_ctx\_regression} $>0$ asks $z$ to predict the
    dataset's $\xi$ via \texttt{ContextRegressionHead}, which fights
    the latent's role as an online posterior; the default is $0$.
\end{itemize}

\section{Implementation pointers}

\paragraph{File layout.}  The folder
\texttt{Robot\_and\_Object/No\_object\_Euler/} contains:
\begin{itemize}
  \item \texttt{models.py} -- new classes \texttt{CouplingHead},
    \texttt{LatentInnovationUpdate}; modified
    \texttt{ContactWrenchHead} (optional context input);
    modified \texttt{RobotObjectMCGDFStep}
    (coupling object dynamics in
    \texttt{coupling\_object\_dynamics}, new \texttt{forward});
    modified \texttt{MultiStepRobotObjectMCGDFWorldModel}
    (\texttt{step\_filter} method, filter-window loop in
    \texttt{forward}).
  \item \texttt{wm\_train.py} -- new CLI flags
    \texttt{--omit\_object\_context} (default on),
    \texttt{--coupling\_*}, \texttt{--init\_rho\_*},
    \texttt{--filter\_steps}, \texttt{--latent\_update\_*}.  The
    rest of the training loop is structurally unchanged.
  \item \texttt{plot\_multistep\_trajectory.py},
    \texttt{animate\_episode.py} -- the checkpoint loader honors the
    new flags; the rollout helper takes a \texttt{filter\_steps}
    argument and runs the closed-loop filter window before the
    open-loop rollout.  A new CLI flag
    \texttt{--filter\_steps} (default $-1$ = use the checkpoint's
    value) lets the practitioner override at eval time.
  \item \texttt{dataset.py} -- unchanged.  The existing layout
    (\texttt{history\_states}, \texttt{history\_torques},
    \texttt{future\_torques}, \texttt{future\_states},
    \texttt{object\_context}) is exactly what the new
    \texttt{MultiStep.forward} needs; the filter consumes the last
    $L_{\mathrm{filter}}$ transitions of the history window.
\end{itemize}

\paragraph{Constructor flags (new).}  In addition to the MCGDF
flags, \texttt{MultiStepRobotObjectMCGDFWorldModel} accepts:
\begin{itemize}
  \item \texttt{omit\_object\_context} (bool, default
    \texttt{True}): strip $\xi$ from the contact head's input.
  \item \texttt{coupling\_hidden\_dim}, \texttt{coupling\_depth}: the
    \texttt{CouplingHead} MLP size (defaults 64, 2).
  \item \texttt{coupling\_components} (1 or 2): single
    rigid-coupling scalar vs.\ two-component
    $(\alpha^{\mathrm{rigid}}, \alpha^{\mathrm{slide}})$.
  \item \texttt{init\_rho\_v}, \texttt{init\_rho\_omega} (default
    $0.95$): initial passive decays.
  \item \texttt{filter\_steps} (int, default $0$): $L_{\mathrm{filter}}$
    for the closed-loop filter window.  Requires
    \texttt{use\_context\_encoder=True}.
  \item \texttt{latent\_update\_hidden\_dim},
    \texttt{latent\_update\_depth}: \texttt{LatentInnovationUpdate}
    cell size (defaults 64, 1).
\end{itemize}

\section{Caveats}

\begin{itemize}
  \item \emph{No object-side Newton--Euler.}  The new integrator
    cannot represent ballistic motion, free spinning after release,
    or motions driven by external impulses other than the gripper.
    For the Franka Lift task this is acceptable; flag explicitly
    before extending to other tasks.
  \item \emph{$\alpha_t$ conflates regimes.}  With
    \texttt{coupling\_components=1}, a single scalar gates both
    rigid carry and sliding push.  Sliding tangential motion under
    contact pressure looks (to the gate) similar to rigid carry,
    which may force the model into a compromise.  The 2-component
    variant decouples these.
  \item \emph{Filter requires
    \texttt{history\_len} $\geq L_{\mathrm{filter}} + 1$.}  The
    filter consumes the last $L_{\mathrm{filter}}$ transitions of
    the history window; the rest is the encoder's warm-up.  An
    error is raised at runtime if the configuration is inconsistent.
  \item \emph{Truncated BPTT length.}  Backpropagating through
    $L_{\mathrm{filter}}$ filter steps and $H$ rollout steps grows
    the autograd graph linearly in $L_{\mathrm{filter}}+H$.  Large
    filter windows on long horizons may need gradient checkpointing.
  \item \emph{Pre-contact baseline assumes a single contact onset
    per episode.}  Inherited from MCGDF; not changed here.
  \item \emph{Robot-side Jacobian still uses the gripper midpoint.}
    Inherited from MCGDF.  In the coupling block, the displacement
    $\mathbf{d}_t = p_{\mathrm{ee}} - p^o_t$ is measured at the
    same midpoint; per-finger contacts would benefit from the
    per-finger Jacobian extension.
\end{itemize}

\section{Relation to the project proposal (Pi-CDD)}

The latent online update Eq.~\eqref{eq:noe-latent-update} is exactly
the time-evolving context channel sketched in Pi-CDD year-1
(Figure 1 of the proposal), implemented as an innovation
update on $z$ rather than as a re-encoding of a sliding history
window.  The coupling-based object dynamics
Eqs.~\eqref{eq:noe-vobj}--\eqref{eq:noe-robj} fit the proposal's
``Action--Reaction'' framing: $\alpha_t$ \emph{is} the model's
quantitative read of how strongly the gripper action couples to the
object reaction at the current step, with the latent $z$ informing
the threshold.

\end{document}
